\documentclass{article}
\documentclass{article}
\usepackage{amsmath,amssymb}
\usepackage{amsmath,amssymb}
\begin{document}
\documentclass[11pt]{article}
\usepackage{amsmath,amssymb,amsthm}

\newtheorem{lemma}{Lemma}

\begin{document}

\begin{lemma}[Unique Extension of Rooted Tree Isomorphism]
Let $(T_1,r_1),(T_2,r_2)$ be rooted trees of maximum degree $\Delta$.
Let $\varphi_0:\{r_1\}\to\{r_2\}$ be the root map.
Assume for every vertex $v\in T_1$ the multiset of isomorphism types of rooted subtrees at children of $v$
coincides with that of $\varphi(v)$ in $T_2$.
Then there exists a unique isomorphism $\varphi:T_1\to T_2$ extending $\varphi_0$.
\end{lemma}

\begin{proof}
Define $\varphi$ inductively by depth. At depth $0$, $\varphi(r_1)=r_2$.
Assume $\varphi$ defined on all vertices at depth $\le d$.
Let $v$ be at depth $d$, $\varphi(v)=w$.
Let $C(v),C(w)$ be children sets. By hypothesis, there exists a bijection
$\psi_v:C(v)\to C(w)$ preserving rooted subtree types.
For $u\in C(v)$, define $\varphi(u):=\psi_v(u)$.
This defines $\varphi$ on depth $d+1$.
By induction, $\varphi$ exists on all vertices.
Uniqueness follows since $\psi_v$ is uniquely determined at each step.
\end{proof}

\begin{lemma}[Transitivity of CFI Gadget Automorphisms]
Let $G_{\mathrm{CFI}}(v)$ be the CFI gadget at vertex $v$ of degree $d$.
Let $S_v$ be the set of local states (valid parity assignments).
Then $\mathrm{Aut}(G_{\mathrm{CFI}}(v))$ acts transitively on $S_v$.
\end{lemma}

\begin{proof}
Model $S_v=\{x\in\mathbb{F}_2^d:\sum x_i=0\}$.
For $i\neq j$, define $\tau_{ij}$ by $\tau_{ij}(x)=x+e_i+e_j$.
These preserve parity, hence lie in $\mathrm{Aut}$.
They generate translations $x\mapsto x+y$ for all $y\in S_v$.
Thus any $x'$ equals $x+y$ for some $y$, proving transitivity.
\end{proof}

\end{document}

\begin{document}


\section*{EF–CFI Transfer Formalization}
\section*{EF–CFI Transfer Formalization}


\textbf{State Space:}
\textbf{State Space:}
\[
\[
S_t = (\bar a_t,\bar b_t),\quad
S_t = (\bar a_t,\bar b_t),\quad
\bar a_t\in V(G_0)^k,\;\bar b_t\in V(G_1)^k,
\bar a_t\in V(G_0)^k,\;\bar b_t\in V(G_1)^k,
\]
\]
with partial isomorphism invariant:
with partial isomorphism invariant:
\[
\[
\forall i,j:\;
\forall i,j:\;
(a_i=a_j)\Leftrightarrow(b_i=b_j),\quad
(a_i=a_j)\Leftrightarrow(b_i=b_j),\quad
E(a_i,a_j)\Leftrightarrow E(b_i,b_j).
E(a_i,a_j)\Leftrightarrow E(b_i,b_j).
\]
\]


\textbf{Transition Function:}
\textbf{Transition Function:}
\[
\[
T:\;S_t\times (V(G_0)\cup V(G_1))\to S_{t+1}
T:\;S_t\times (V(G_0)\cup V(G_1))\to S_{t+1}
\]
\]
defined by extending tuples and truncating to last \(k\) entries.
defined by extending tuples and truncating to last \(k\) entries.


\textbf{Extension Lemma (Tree Neighborhoods):}
\textbf{Extension Lemma (Tree Neighborhoods):}
\[
\[
\mathrm{girth}(B)>2r\Rightarrow
\mathrm{girth}(B)>2r\Rightarrow
B_r(v)\cong T_r(v)
B_r(v)\cong T_r(v)
\Rightarrow
\Rightarrow
\exists!\;\text{isomorphism extension preserving adjacency}.
\exists!\;\text{isomorphism extension preserving adjacency}.
\]
\]


\textbf{Duplicator Response:}
\textbf{Duplicator Response:}
\[
\[
\forall \text{ Spoiler move }x\in G_0,\;
\forall \text{ Spoiler move }x\in G_0,\;
\exists y\in G_1:
\exists y\in G_1:
\]
\]
\[
\[
B_r(x)\cong B_r(y)
B_r(x)\cong B_r(y)
\;\land\;
\;\land\;
S_{t+1}\text{ remains partial isomorphism}.
S_{t+1}\text{ remains partial isomorphism}.
\]
\]


\textbf{CFI Symmetry Invariance:}
\textbf{CFI Symmetry Invariance:}
\[
\[
\text{local automorphisms act transitively on gadget states}
\text{local automorphisms act transitively on gadget states}
\Rightarrow
\Rightarrow
\text{parity indistinguishable locally}.
\text{parity indistinguishable locally}.
\]
\]


\textbf{Induction:}
\textbf{Induction:}
\[
\[
(\forall t\le k)\;S_t\text{ is partial isomorphism}
(\forall t\le k)\;S_t\text{ is partial isomorphism}
\Rightarrow
\Rightarrow
\text{Duplicator wins }k\text{-round EF game}.
\text{Duplicator wins }k\text{-round EF game}.
\]
\]


\textbf{Conclusion:}
\textbf{Conclusion:}
\[
\[
G_0\equiv_{FO^k}G_1.
G_0\equiv_{FO^k}G_1.
\]
\]


\textbf{Status: OPEN (requires formal proof of uniqueness and automorphism action).}
\textbf{Status: OPEN (requires formal proof of uniqueness and automorphism action).}


\begin{theorem}[EF--CFI Separation]
Let G_0, G_1 be CFI lifts over a common base graph.
Then G_0 \equiv_{FO^k} G_1 and I(G_0) \neq I(G_1).
\end{theorem}

\begin{proof}
FO^k equivalence follows from the EF-game strategy using Lemma~\ref{lem:tree-extension} and Lemma~\ref{lem:cfi-transitivity}.
Invariant separation follows from the nontrivial cohomology witness.
\end{proof}
\end{document}
\end{document}
